%%%%%%%%%%%%%%%%%%%%%%%%%%%%%%%%%%%%%%%%%%%%%%%%%%%%%%%%%%%%
% 5.18 GEOMETRIC ANALYSIS
% The Composition of Advanced Mathematics
%%%%%%%%%%%%%%%%%%%%%%%%%%%%%%%%%%%%%%%%%%%%%%%%%%%%%%%%%%%%

\section{Geometric Analysis}
\label{sec:geometric-analysis}

\prerequisite{
Multivariable calculus, vector calculus, real analysis, differential
geometry, manifolds, functional analysis, calculus of variations, and
basic partial differential equations.
}

\learningobjectives{
By the end of this section, the reader should be able to:
\begin{itemize}
    \item explain the central goals of geometric analysis;
    \item understand how geometric problems lead to differential equations;
    \item work with Riemannian metrics, volume forms, and geometric operators;
    \item understand the gradient, divergence, and Laplace--Beltrami operator;
    \item formulate Dirichlet energy on a Riemannian manifold;
    \item recognize harmonic functions and harmonic maps;
    \item understand geodesics from a variational viewpoint;
    \item describe minimal surfaces and mean curvature;
    \item understand the first variation of area;
    \item recognize elliptic equations arising from geometry;
    \item understand curvature as an analytic quantity;
    \item distinguish sectional, Ricci, and scalar curvature;
    \item understand the basic structure of the Yamabe problem;
    \item describe mean curvature flow and Ricci flow;
    \item understand geometric evolution equations;
    \item recognize maximum principles and energy methods;
    \item understand weak solutions and regularity questions;
    \item recognize concentration and compactness phenomena;
    \item understand geometric inequalities and comparison principles;
    \item connect geometric analysis with topology, PDEs,
          mathematical physics, and differential geometry.
\end{itemize}
}

%%%%%%%%%%%%%%%%%%%%%%%%%%%%%%%%%%%%%%%%%%%%%%%%%%%%%%%%%%%%
\subsection{What Is Geometric Analysis?}
\label{subsec:what-is-geometric-analysis}
%%%%%%%%%%%%%%%%%%%%%%%%%%%%%%%%%%%%%%%%%%%%%%%%%%%%%%%%%%%%

Geometric analysis studies geometric objects using the methods of
analysis.

Its central philosophy is

\[
\boxed{
\text{geometry}
\longleftrightarrow
\text{differential equations}.
}
\]

A geometric condition is translated into an analytic equation.

The equation is then studied using tools such as

\[
\boxed{
\text{calculus of variations},
\quad
\text{PDE theory},
\quad
\text{functional analysis},
\quad
\text{measure theory}.
}
\]

The resulting analytic information is interpreted geometrically.

%%%%%%%%%%%%%%%%%%%%%%%%%%%%%%%%%%%%%%%%%%%%%%%%%%%%%%%%%%%%
\subsection{Typical Problems}
\label{subsec:geometric-analysis-problems}
%%%%%%%%%%%%%%%%%%%%%%%%%%%%%%%%%%%%%%%%%%%%%%%%%%%%%%%%%%%%

Important problems in geometric analysis include:

\begin{itemize}
    \item finding shortest curves;
    \item finding surfaces of least area;
    \item constructing metrics with prescribed curvature;
    \item studying harmonic functions on manifolds;
    \item studying harmonic maps between manifolds;
    \item evolving geometric structures to improve their shape;
    \item understanding singularities of geometric flows;
    \item relating curvature to topology.
\end{itemize}

These problems often become nonlinear differential equations.

%%%%%%%%%%%%%%%%%%%%%%%%%%%%%%%%%%%%%%%%%%%%%%%%%%%%%%%%%%%%
\subsection{Riemannian Manifolds}
\label{subsec:geometric-analysis-riemannian-manifolds}
%%%%%%%%%%%%%%%%%%%%%%%%%%%%%%%%%%%%%%%%%%%%%%%%%%%%%%%%%%%%

Let \(M\) be a smooth manifold.

A Riemannian metric assigns an inner product

\[
g_p
\]

to each tangent space

\[
T_pM.
\]

Locally,

\[
\boxed{
g
=
\sum_{i,j}
g_{ij}\,dx^i\,dx^j.
}
\]

The matrix

\[
(g_{ij})
\]

is symmetric and positive definite.

The metric determines

\[
\boxed{
\text{length},
\quad
\text{angle},
\quad
\text{distance},
\quad
\text{volume},
\quad
\text{curvature}.
}
\]

%%%%%%%%%%%%%%%%%%%%%%%%%%%%%%%%%%%%%%%%%%%%%%%%%%%%%%%%%%%%
\subsection{Inverse Metric}
\label{subsec:inverse-metric}
%%%%%%%%%%%%%%%%%%%%%%%%%%%%%%%%%%%%%%%%%%%%%%%%%%%%%%%%%%%%

The inverse matrix of

\[
(g_{ij})
\]

is written

\[
(g^{ij}),
\]

so that

\[
\boxed{
\sum_j
g^{ij}g_{jk}
=
\delta^i_k.
}
\]

Here \(\delta^i_k\), called the Kronecker delta, equals \(1\) when
\(i=k\) and \(0\) otherwise.

The inverse metric is used to convert covariant components into
contravariant components.

%%%%%%%%%%%%%%%%%%%%%%%%%%%%%%%%%%%%%%%%%%%%%%%%%%%%%%%%%%%%
\subsection{Riemannian Volume}
\label{subsec:riemannian-volume}
%%%%%%%%%%%%%%%%%%%%%%%%%%%%%%%%%%%%%%%%%%%%%%%%%%%%%%%%%%%%

The Riemannian metric determines a natural volume element.

In local coordinates,

\[
\boxed{
dV_g
=
\sqrt{\det(g_{ij})}\,
dx^1\cdots dx^n.
}
\]

Thus integration on a Riemannian manifold takes the form

\[
\boxed{
\int_M f\,dV_g.
}
\]

The metric therefore influences not only distances but also integration.

%%%%%%%%%%%%%%%%%%%%%%%%%%%%%%%%%%%%%%%%%%%%%%%%%%%%%%%%%%%%
\subsection{Gradient on a Riemannian Manifold}
\label{subsec:riemannian-gradient}
%%%%%%%%%%%%%%%%%%%%%%%%%%%%%%%%%%%%%%%%%%%%%%%%%%%%%%%%%%%%

For a smooth function

\[
f:M\to\R,
\]

the gradient \(\nabla f\) is defined by

\[
\boxed{
g(\nabla f,X)
=
df(X)
}
\]

for every vector field \(X\).

In local coordinates,

\[
\boxed{
\nabla f
=
\sum_{i,j}
g^{ij}
\frac{\partial f}{\partial x^j}
\frac{\partial}{\partial x^i}.
}
\]

Thus the gradient depends on the metric.

%%%%%%%%%%%%%%%%%%%%%%%%%%%%%%%%%%%%%%%%%%%%%%%%%%%%%%%%%%%%
\subsection{Divergence}
\label{subsec:riemannian-divergence}
%%%%%%%%%%%%%%%%%%%%%%%%%%%%%%%%%%%%%%%%%%%%%%%%%%%%%%%%%%%%

For a vector field

\[
X
=
X^i
\frac{\partial}{\partial x^i},
\]

the Riemannian divergence is locally

\[
\boxed{
\operatorname{div}_g X
=
\frac{1}{\sqrt{\det g}}
\frac{\partial}{\partial x^i}
\left(
\sqrt{\det g}\,X^i
\right),
}
\]

where repeated indices are summed.

The divergence measures the local expansion or contraction of the vector
field relative to the Riemannian volume.

%%%%%%%%%%%%%%%%%%%%%%%%%%%%%%%%%%%%%%%%%%%%%%%%%%%%%%%%%%%%
\subsection{Laplace--Beltrami Operator}
\label{subsec:laplace-beltrami}
%%%%%%%%%%%%%%%%%%%%%%%%%%%%%%%%%%%%%%%%%%%%%%%%%%%%%%%%%%%%

\begin{definitionbox}{Laplace--Beltrami Operator}
For a smooth function \(f\) on a Riemannian manifold,

\[
\boxed{
\Delta_g f
=
\operatorname{div}_g(\nabla f).
}
\]
\end{definitionbox}

In local coordinates,

\[
\boxed{
\Delta_g f
=
\frac{1}{\sqrt{\det g}}
\frac{\partial}{\partial x^i}
\left(
\sqrt{\det g}\,
g^{ij}
\frac{\partial f}{\partial x^j}
\right).
}
\]

This generalizes the ordinary Euclidean Laplacian.

%%%%%%%%%%%%%%%%%%%%%%%%%%%%%%%%%%%%%%%%%%%%%%%%%%%%%%%%%%%%
\subsection{Euclidean Special Case}
\label{subsec:laplace-beltrami-euclidean}
%%%%%%%%%%%%%%%%%%%%%%%%%%%%%%%%%%%%%%%%%%%%%%%%%%%%%%%%%%%%

For Euclidean space,

\[
g_{ij}
=
\delta_{ij}.
\]

Hence

\[
\det g=1
\]

and

\[
g^{ij}=\delta^{ij}.
\]

Therefore,

\[
\boxed{
\Delta_g f
=
\sum_{i=1}^{n}
\frac{\partial^2f}{\partial(x^i)^2}.
}
\]

Thus the Laplace--Beltrami operator reduces to the ordinary Laplacian.

%%%%%%%%%%%%%%%%%%%%%%%%%%%%%%%%%%%%%%%%%%%%%%%%%%%%%%%%%%%%
\subsection{Harmonic Functions}
\label{subsec:geometric-harmonic-functions}
%%%%%%%%%%%%%%%%%%%%%%%%%%%%%%%%%%%%%%%%%%%%%%%%%%%%%%%%%%%%

\begin{definitionbox}{Harmonic Function}
A smooth function \(u\) on a Riemannian manifold is harmonic if

\[
\boxed{
\Delta_g u=0.
}
\]
\end{definitionbox}

Harmonic functions generalize classical solutions of Laplace's equation.

Their behavior is strongly influenced by the geometry of the manifold.

%%%%%%%%%%%%%%%%%%%%%%%%%%%%%%%%%%%%%%%%%%%%%%%%%%%%%%%%%%%%
\subsection{Dirichlet Energy}
\label{subsec:geometric-dirichlet-energy}
%%%%%%%%%%%%%%%%%%%%%%%%%%%%%%%%%%%%%%%%%%%%%%%%%%%%%%%%%%%%

Define the Dirichlet energy by

\[
\boxed{
E[u]
=
\frac12
\int_M
|\nabla u|_g^2
\,dV_g.
}
\]

Critical points of this functional satisfy

\[
\boxed{
\Delta_g u=0
}
\]

under appropriate boundary conditions.

Thus harmonic functions arise naturally from a variational problem.

%%%%%%%%%%%%%%%%%%%%%%%%%%%%%%%%%%%%%%%%%%%%%%%%%%%%%%%%%%%%
\subsection{First Variation of Dirichlet Energy}
\label{subsec:first-variation-dirichlet}
%%%%%%%%%%%%%%%%%%%%%%%%%%%%%%%%%%%%%%%%%%%%%%%%%%%%%%%%%%%%

Consider a variation

\[
u_\varepsilon
=
u+\varepsilon\eta.
\]

Then

\[
E[u_\varepsilon]
=
\frac12
\int_M
|\nabla u+\varepsilon\nabla\eta|_g^2
\,dV_g.
\]

Differentiating at

\[
\varepsilon=0
\]

gives

\[
\boxed{
\delta E[u;\eta]
=
\int_M
g(\nabla u,\nabla\eta)
\,dV_g.
}
\]

Integration by parts gives, for suitable variations,

\[
\boxed{
\delta E[u;\eta]
=
-
\int_M
(\Delta_g u)\eta
\,dV_g.
}
\]

Therefore stationarity implies

\[
\boxed{
\Delta_g u=0.
}
\]

%%%%%%%%%%%%%%%%%%%%%%%%%%%%%%%%%%%%%%%%%%%%%%%%%%%%%%%%%%%%
\subsection{Geodesics}
\label{subsec:geometric-analysis-geodesics}
%%%%%%%%%%%%%%%%%%%%%%%%%%%%%%%%%%%%%%%%%%%%%%%%%%%%%%%%%%%%

A geodesic is the Riemannian analogue of a straight line.

For a curve

\[
\gamma:[a,b]\to M,
\]

its length is

\[
\boxed{
L[\gamma]
=
\int_a^b
|\dot{\gamma}(t)|_g
\,dt.
}
\]

Its energy is

\[
\boxed{
E[\gamma]
=
\frac12
\int_a^b
|\dot{\gamma}(t)|_g^2
\,dt.
}
\]

Critical points of the energy satisfy the geodesic equation.

%%%%%%%%%%%%%%%%%%%%%%%%%%%%%%%%%%%%%%%%%%%%%%%%%%%%%%%%%%%%
\subsection{Geodesic Equation}
\label{subsec:geodesic-equation-analysis}
%%%%%%%%%%%%%%%%%%%%%%%%%%%%%%%%%%%%%%%%%%%%%%%%%%%%%%%%%%%%

In local coordinates,

\[
\gamma(t)
=
\left(
x^1(t),\ldots,x^n(t)
\right).
\]

The geodesic equation is

\[
\boxed{
\frac{d^2x^k}{dt^2}
+
\Gamma^k_{ij}
\frac{dx^i}{dt}
\frac{dx^j}{dt}
=
0.
}
\]

The quantities

\[
\Gamma^k_{ij}
\]

are the Christoffel symbols of the Levi-Civita connection.

They are determined by the metric.

%%%%%%%%%%%%%%%%%%%%%%%%%%%%%%%%%%%%%%%%%%%%%%%%%%%%%%%%%%%%
\subsection{Christoffel Symbols}
\label{subsec:christoffel-symbols-analysis}
%%%%%%%%%%%%%%%%%%%%%%%%%%%%%%%%%%%%%%%%%%%%%%%%%%%%%%%%%%%%

The Christoffel symbols are

\[
\boxed{
\Gamma^k_{ij}
=
\frac12
g^{k\ell}
\left(
\frac{\partial g_{j\ell}}{\partial x^i}
+
\frac{\partial g_{i\ell}}{\partial x^j}
-
\frac{\partial g_{ij}}{\partial x^\ell}
\right).
}
\]

They are not tensor components by themselves.

Instead, they encode the Levi-Civita connection in a coordinate system.

%%%%%%%%%%%%%%%%%%%%%%%%%%%%%%%%%%%%%%%%%%%%%%%%%%%%%%%%%%%%
\subsection{Minimal Surfaces}
\label{subsec:geometric-minimal-surfaces}
%%%%%%%%%%%%%%%%%%%%%%%%%%%%%%%%%%%%%%%%%%%%%%%%%%%%%%%%%%%%

Minimal surfaces are critical points of the area functional.

For a graph

\[
z=u(x,y),
\]

the area is

\[
\boxed{
A[u]
=
\int_\Omega
\sqrt{
1+u_x^2+u_y^2
}
\,dx\,dy.
}
\]

The Euler--Lagrange equation is

\[
\boxed{
\operatorname{div}
\left(
\frac{\nabla u}
{\sqrt{1+|\nabla u|^2}}
\right)
=
0.
}
\]

This is the minimal surface equation.

%%%%%%%%%%%%%%%%%%%%%%%%%%%%%%%%%%%%%%%%%%%%%%%%%%%%%%%%%%%%
\subsection{Mean Curvature}
\label{subsec:mean-curvature}
%%%%%%%%%%%%%%%%%%%%%%%%%%%%%%%%%%%%%%%%%%%%%%%%%%%%%%%%%%%%

Mean curvature measures how a hypersurface bends inside its ambient
space.

If

\[
\kappa_1,\ldots,\kappa_m
\]

are the principal curvatures, one common convention defines

\[
\boxed{
H
=
\frac{1}{m}
\sum_{j=1}^{m}\kappa_j.
}
\]

Some literature instead uses the unnormalized sum.

Therefore conventions must always be checked.

A minimal hypersurface satisfies

\[
\boxed{
H=0.
}
\]

%%%%%%%%%%%%%%%%%%%%%%%%%%%%%%%%%%%%%%%%%%%%%%%%%%%%%%%%%%%%
\subsection{First Variation of Area}
\label{subsec:first-variation-area}
%%%%%%%%%%%%%%%%%%%%%%%%%%%%%%%%%%%%%%%%%%%%%%%%%%%%%%%%%%%%

Let \(\Sigma\) be a hypersurface and consider a normal variation with
normal speed \(\phi\).

Under a standard sign convention, the first variation of area has the
form

\[
\boxed{
\left.
\frac{d}{dt}
\operatorname{Area}(\Sigma_t)
\right|_{t=0}
=
-
\int_\Sigma
H_{\mathrm{sum}}\phi
\,d\mu,
}
\]

where \(H_{\mathrm{sum}}\) denotes the unnormalized mean curvature
associated with the chosen unit normal.

Therefore a critical point of area satisfies

\[
\boxed{
H_{\mathrm{sum}}=0.
}
\]

This is the geometric Euler--Lagrange equation for minimal hypersurfaces.

%%%%%%%%%%%%%%%%%%%%%%%%%%%%%%%%%%%%%%%%%%%%%%%%%%%%%%%%%%%%
\subsection{Worked Example: Plane as a Minimal Surface}
\label{subsec:worked-plane-minimal}
%%%%%%%%%%%%%%%%%%%%%%%%%%%%%%%%%%%%%%%%%%%%%%%%%%%%%%%%%%%%

\begin{workedexamplebox}{A Plane Has Zero Mean Curvature}

Consider

\[
u(x,y)=ax+by+c.
\]

Then

\[
\nabla u=(a,b)
\]

is constant.

Therefore,

\[
\frac{\nabla u}
{\sqrt{1+|\nabla u|^2}}
\]

is also constant.

Its divergence is zero.

Hence the minimal surface equation is satisfied.

\begin{answerbox}
\[
\boxed{
\text{Every plane is a minimal surface.}
}
\]
\end{answerbox}

\end{workedexamplebox}

%%%%%%%%%%%%%%%%%%%%%%%%%%%%%%%%%%%%%%%%%%%%%%%%%%%%%%%%%%%%
\subsection{Harmonic Maps}
\label{subsec:harmonic-maps}
%%%%%%%%%%%%%%%%%%%%%%%%%%%%%%%%%%%%%%%%%%%%%%%%%%%%%%%%%%%%

Let

\[
(M,g)
\]

and

\[
(N,h)
\]

be Riemannian manifolds.

A smooth map

\[
u:M\to N
\]

has energy

\[
\boxed{
E[u]
=
\frac12
\int_M
|du|^2
\,dV_g.
}
\]

Critical points of this energy are called harmonic maps.

They generalize harmonic functions.

%%%%%%%%%%%%%%%%%%%%%%%%%%%%%%%%%%%%%%%%%%%%%%%%%%%%%%%%%%%%
\subsection{Tension Field}
\label{subsec:tension-field}
%%%%%%%%%%%%%%%%%%%%%%%%%%%%%%%%%%%%%%%%%%%%%%%%%%%%%%%%%%%%

The Euler--Lagrange equation for harmonic maps is written

\[
\boxed{
\tau(u)=0,
}
\]

where

\[
\tau(u)
\]

is called the tension field.

In local coordinates, the equation has the schematic form

\[
\boxed{
\Delta_g u^\alpha
+
\Gamma^\alpha_{\beta\gamma}(u)
g^{ij}
\frac{\partial u^\beta}{\partial x^i}
\frac{\partial u^\gamma}{\partial x^j}
=
0.
}
\]

Thus harmonic-map equations are generally nonlinear.

%%%%%%%%%%%%%%%%%%%%%%%%%%%%%%%%%%%%%%%%%%%%%%%%%%%%%%%%%%%%
\subsection{Curvature}
\label{subsec:geometric-analysis-curvature}
%%%%%%%%%%%%%%%%%%%%%%%%%%%%%%%%%%%%%%%%%%%%%%%%%%%%%%%%%%%%

Curvature measures the failure of a geometric space to behave locally
like Euclidean space beyond first order.

The Riemann curvature tensor is written schematically as

\[
\boxed{
\operatorname{Rm}.
}
\]

It contains the full local curvature information of a Riemannian
manifold.

Several important contractions and derived quantities arise from it.

%%%%%%%%%%%%%%%%%%%%%%%%%%%%%%%%%%%%%%%%%%%%%%%%%%%%%%%%%%%%
\subsection{Sectional Curvature}
\label{subsec:sectional-curvature-analysis}
%%%%%%%%%%%%%%%%%%%%%%%%%%%%%%%%%%%%%%%%%%%%%%%%%%%%%%%%%%%%

For a two-dimensional tangent plane

\[
\sigma\subseteq T_pM,
\]

the sectional curvature is

\[
\boxed{
K(\sigma).
}
\]

It measures curvature in the direction of the chosen tangent plane.

In dimension two, sectional curvature reduces to Gaussian curvature.

%%%%%%%%%%%%%%%%%%%%%%%%%%%%%%%%%%%%%%%%%%%%%%%%%%%%%%%%%%%%
\subsection{Ricci Curvature}
\label{subsec:ricci-curvature-analysis}
%%%%%%%%%%%%%%%%%%%%%%%%%%%%%%%%%%%%%%%%%%%%%%%%%%%%%%%%%%%%

The Ricci tensor is obtained by tracing the Riemann curvature tensor.

It is written

\[
\boxed{
\operatorname{Ric}.
}
\]

Ricci curvature measures how volumes and families of nearby geodesics
deviate from their Euclidean behavior.

It plays a central role in

\[
\boxed{
\text{comparison geometry},
\quad
\text{general relativity},
\quad
\text{Ricci flow}.
}
\]

%%%%%%%%%%%%%%%%%%%%%%%%%%%%%%%%%%%%%%%%%%%%%%%%%%%%%%%%%%%%
\subsection{Scalar Curvature}
\label{subsec:scalar-curvature-analysis}
%%%%%%%%%%%%%%%%%%%%%%%%%%%%%%%%%%%%%%%%%%%%%%%%%%%%%%%%%%%%

Scalar curvature is the trace of the Ricci tensor:

\[
\boxed{
R
=
g^{ij}\operatorname{Ric}_{ij}.
}
\]

It is a scalar function on the manifold.

The hierarchy is therefore

\[
\boxed{
\operatorname{Rm}
\longrightarrow
\operatorname{Ric}
\longrightarrow
R.
}
\]

Each contraction retains less information but often leads to more
manageable equations.

%%%%%%%%%%%%%%%%%%%%%%%%%%%%%%%%%%%%%%%%%%%%%%%%%%%%%%%%%%%%
\subsection{Curvature and Differential Equations}
\label{subsec:curvature-differential-equations}
%%%%%%%%%%%%%%%%%%%%%%%%%%%%%%%%%%%%%%%%%%%%%%%%%%%%%%%%%%%%

Many geometric-analysis problems seek a metric \(g\) satisfying a
curvature equation.

Examples include

\[
\boxed{
\operatorname{Ric}(g)=0,
}
\]

\[
\boxed{
\operatorname{Ric}(g)=\lambda g,
}
\]

and

\[
\boxed{
R(g)=\text{prescribed function}.
}
\]

These are nonlinear differential equations for the components of the
metric.

%%%%%%%%%%%%%%%%%%%%%%%%%%%%%%%%%%%%%%%%%%%%%%%%%%%%%%%%%%%%
\subsection{Einstein Metrics}
\label{subsec:einstein-metrics}
%%%%%%%%%%%%%%%%%%%%%%%%%%%%%%%%%%%%%%%%%%%%%%%%%%%%%%%%%%%%

\begin{definitionbox}{Einstein Metric}
A Riemannian metric is Einstein if

\[
\boxed{
\operatorname{Ric}
=
\lambda g
}
\]

for some constant \(\lambda\).
\end{definitionbox}

Einstein metrics are important in differential geometry and mathematical
physics.

The equation is nonlinear because the Ricci tensor itself depends on the
metric and its derivatives.

%%%%%%%%%%%%%%%%%%%%%%%%%%%%%%%%%%%%%%%%%%%%%%%%%%%%%%%%%%%%
\subsection{Conformal Changes of Metric}
\label{subsec:conformal-change-metric}
%%%%%%%%%%%%%%%%%%%%%%%%%%%%%%%%%%%%%%%%%%%%%%%%%%%%%%%%%%%%

Two metrics are conformally related if

\[
\boxed{
\widetilde g
=
e^{2u}g
}
\]

for some smooth function \(u\).

Equivalently, in dimensions where it is convenient, one often writes

\[
\boxed{
\widetilde g
=
v^{\frac{4}{n-2}}g,
\qquad
n\geq3.
}
\]

Conformal changes preserve angles but alter lengths and curvature.

This converts geometric curvature questions into nonlinear PDEs.

%%%%%%%%%%%%%%%%%%%%%%%%%%%%%%%%%%%%%%%%%%%%%%%%%%%%%%%%%%%%
\subsection{The Yamabe Problem}
\label{subsec:yamabe-problem}
%%%%%%%%%%%%%%%%%%%%%%%%%%%%%%%%%%%%%%%%%%%%%%%%%%%%%%%%%%%%

The Yamabe problem asks whether every compact Riemannian manifold of
dimension

\[
n\geq3
\]

admits a conformally equivalent metric with constant scalar curvature.

Writing

\[
\widetilde g
=
u^{\frac{4}{n-2}}g,
\qquad
u>0,
\]

leads to an equation of the form

\[
\boxed{
-c_n\Delta_g u
+
R_g u
=
\lambda
u^{\frac{n+2}{n-2}},
}
\]

where

\[
\boxed{
c_n
=
\frac{4(n-1)}{n-2}.
}
\]

This is a nonlinear elliptic equation involving the critical Sobolev
exponent.

%%%%%%%%%%%%%%%%%%%%%%%%%%%%%%%%%%%%%%%%%%%%%%%%%%%%%%%%%%%%
\subsection{Yamabe Functional}
\label{subsec:yamabe-functional}
%%%%%%%%%%%%%%%%%%%%%%%%%%%%%%%%%%%%%%%%%%%%%%%%%%%%%%%%%%%%

The Yamabe equation arises variationally from a quotient of the form

\[
\boxed{
Q_g[u]
=
\frac{
\displaystyle
\int_M
\left(
c_n|\nabla u|^2
+
R_gu^2
\right)
\,dV_g
}{
\displaystyle
\left(
\int_M
|u|^{\frac{2n}{n-2}}
\,dV_g
\right)^{\frac{n-2}{n}}
}.
}
\]

This example illustrates a recurring pattern:

\[
\boxed{
\text{geometric condition}
\longrightarrow
\text{variational problem}
\longrightarrow
\text{nonlinear PDE}.
}
\]

%%%%%%%%%%%%%%%%%%%%%%%%%%%%%%%%%%%%%%%%%%%%%%%%%%%%%%%%%%%%
\subsection{Elliptic Equations in Geometry}
\label{subsec:elliptic-equations-geometry}
%%%%%%%%%%%%%%%%%%%%%%%%%%%%%%%%%%%%%%%%%%%%%%%%%%%%%%%%%%%%

Many geometric equations are elliptic after suitable coordinate choices
or gauge conditions.

Elliptic equations tend to have strong regularity properties.

Typical analytic questions include:

\[
\boxed{
\text{Does a solution exist?}
}
\]

\[
\boxed{
\text{Is it unique?}
}
\]

\[
\boxed{
\text{How smooth is it?}
}
\]

\[
\boxed{
\text{Can its size be estimated?}
}
\]

These analytic questions often determine geometric structure.

%%%%%%%%%%%%%%%%%%%%%%%%%%%%%%%%%%%%%%%%%%%%%%%%%%%%%%%%%%%%
\subsection{Maximum Principles}
\label{subsec:geometric-maximum-principles}
%%%%%%%%%%%%%%%%%%%%%%%%%%%%%%%%%%%%%%%%%%%%%%%%%%%%%%%%%%%%

Maximum principles are among the most important tools in geometric
analysis.

For a harmonic function satisfying

\[
\Delta_g u=0
\]

on a suitable connected domain, interior maxima and minima are strongly
restricted.

A typical maximum principle implies

\[
\boxed{
\max_{\overline{\Omega}}u
=
\max_{\partial\Omega}u
}
\]

under suitable assumptions.

Such results can imply uniqueness and geometric rigidity.

%%%%%%%%%%%%%%%%%%%%%%%%%%%%%%%%%%%%%%%%%%%%%%%%%%%%%%%%%%%%
\subsection{Bochner Formula}
\label{subsec:bochner-formula}
%%%%%%%%%%%%%%%%%%%%%%%%%%%%%%%%%%%%%%%%%%%%%%%%%%%%%%%%%%%%

A fundamental identity in geometric analysis is the Bochner formula.

For a smooth function \(u\),

\[
\boxed{
\frac12
\Delta_g|\nabla u|^2
=
|\operatorname{Hess}u|^2
+
g(\nabla u,\nabla\Delta_g u)
+
\operatorname{Ric}(\nabla u,\nabla u).
}
\]

Here

\[
\operatorname{Hess}u
\]

denotes the Hessian of \(u\).

This identity connects

\[
\boxed{
\text{second derivatives}
+
\text{curvature}
+
\text{Laplacian}.
}
\]

%%%%%%%%%%%%%%%%%%%%%%%%%%%%%%%%%%%%%%%%%%%%%%%%%%%%%%%%%%%%
\subsection{Why the Bochner Formula Matters}
\label{subsec:why-bochner-formula-matters}
%%%%%%%%%%%%%%%%%%%%%%%%%%%%%%%%%%%%%%%%%%%%%%%%%%%%%%%%%%%%

Suppose \(u\) is harmonic.

Then

\[
\Delta_g u=0,
\]

so the Bochner formula becomes

\[
\boxed{
\frac12
\Delta_g|\nabla u|^2
=
|\operatorname{Hess}u|^2
+
\operatorname{Ric}(\nabla u,\nabla u).
}
\]

If

\[
\operatorname{Ric}\geq0,
\]

the right-hand side is nonnegative.

Thus curvature information produces analytic inequalities.

This is a central mechanism in geometric analysis.

%%%%%%%%%%%%%%%%%%%%%%%%%%%%%%%%%%%%%%%%%%%%%%%%%%%%%%%%%%%%
\subsection{Geometric Inequalities}
\label{subsec:geometric-inequalities}
%%%%%%%%%%%%%%%%%%%%%%%%%%%%%%%%%%%%%%%%%%%%%%%%%%%%%%%%%%%%

Geometric analysis contains many inequalities relating geometric
quantities.

Examples include inequalities involving

\[
\boxed{
\text{area},
\quad
\text{volume},
\quad
\text{curvature},
\quad
\text{eigenvalues},
\quad
\text{diameter}.
}
\]

A classical Euclidean example is the isoperimetric inequality

\[
\boxed{
A
\leq
\frac{L^2}{4\pi}
}
\]

for a planar region with area \(A\) and boundary length \(L\).

Equality occurs for a circle.

%%%%%%%%%%%%%%%%%%%%%%%%%%%%%%%%%%%%%%%%%%%%%%%%%%%%%%%%%%%%
\subsection{Sobolev Inequalities on Manifolds}
\label{subsec:sobolev-inequalities-manifolds}
%%%%%%%%%%%%%%%%%%%%%%%%%%%%%%%%%%%%%%%%%%%%%%%%%%%%%%%%%%%%

Sobolev inequalities relate the size of a function to the size of its
derivatives.

A representative Euclidean form for

\[
n>2
\]

is

\[
\boxed{
\|u\|_{L^{\frac{2n}{n-2}}}
\leq
C
\|\nabla u\|_{L^2}.
}
\]

Analogous inequalities on manifolds depend on geometric information.

They are essential in nonlinear geometric PDEs.

%%%%%%%%%%%%%%%%%%%%%%%%%%%%%%%%%%%%%%%%%%%%%%%%%%%%%%%%%%%%
\subsection{Poincar\'e Inequalities}
\label{subsec:poincare-inequalities-geometric}
%%%%%%%%%%%%%%%%%%%%%%%%%%%%%%%%%%%%%%%%%%%%%%%%%%%%%%%%%%%%

A Poincar\'e inequality controls a function by its derivatives after
removing an appropriate constant component.

A typical form is

\[
\boxed{
\|u-u_M\|_{L^2(M)}
\leq
C
\|\nabla u\|_{L^2(M)},
}
\]

where \(u_M\) is the average value of \(u\).

Such inequalities connect geometry, spectral theory, and PDE estimates.

%%%%%%%%%%%%%%%%%%%%%%%%%%%%%%%%%%%%%%%%%%%%%%%%%%%%%%%%%%%%
\subsection{Spectrum of the Laplace--Beltrami Operator}
\label{subsec:laplace-beltrami-spectrum}
%%%%%%%%%%%%%%%%%%%%%%%%%%%%%%%%%%%%%%%%%%%%%%%%%%%%%%%%%%%%

On a compact manifold, one studies eigenvalue equations of the form

\[
\boxed{
-\Delta_g\phi
=
\lambda\phi.
}
\]

Under standard assumptions, the eigenvalues form a sequence

\[
\boxed{
0\leq
\lambda_0
\leq
\lambda_1
\leq
\lambda_2
\leq
\cdots
}
\]

with

\[
\lambda_k\to\infty.
\]

The spectrum contains geometric information about the manifold.

%%%%%%%%%%%%%%%%%%%%%%%%%%%%%%%%%%%%%%%%%%%%%%%%%%%%%%%%%%%%
\subsection{Rayleigh Quotient}
\label{subsec:geometric-rayleigh-quotient}
%%%%%%%%%%%%%%%%%%%%%%%%%%%%%%%%%%%%%%%%%%%%%%%%%%%%%%%%%%%%

For a nonzero function \(u\), define

\[
\boxed{
R[u]
=
\frac{
\int_M|\nabla u|^2\,dV_g
}{
\int_M|u|^2\,dV_g
}.
}
\]

Variational principles characterize Laplace eigenvalues using this
quotient.

For example, after excluding constant functions on a connected compact
manifold without boundary, the first positive eigenvalue can be
characterized by minimizing \(R[u]\) over functions orthogonal to the
constants.

%%%%%%%%%%%%%%%%%%%%%%%%%%%%%%%%%%%%%%%%%%%%%%%%%%%%%%%%%%%%
\subsection{Geometry and Eigenvalues}
\label{subsec:geometry-eigenvalues}
%%%%%%%%%%%%%%%%%%%%%%%%%%%%%%%%%%%%%%%%%%%%%%%%%%%%%%%%%%%%

Eigenvalues of geometric operators depend on properties such as

\[
\boxed{
\text{volume},
\quad
\text{diameter},
\quad
\text{curvature},
\quad
\text{topology}.
}
\]

Conversely, spectral information can reveal aspects of geometry.

This interaction is the subject of spectral geometry.

%%%%%%%%%%%%%%%%%%%%%%%%%%%%%%%%%%%%%%%%%%%%%%%%%%%%%%%%%%%%
\subsection{Geometric Flows}
\label{subsec:geometric-flows}
%%%%%%%%%%%%%%%%%%%%%%%%%%%%%%%%%%%%%%%%%%%%%%%%%%%%%%%%%%%%

A geometric flow evolves a geometric object according to a differential
equation.

The goal is often to smooth irregularities or drive the object toward a
canonical form.

Two fundamental examples are

\[
\boxed{
\text{mean curvature flow}
}
\]

and

\[
\boxed{
\text{Ricci flow}.
}
\]

These are nonlinear evolution equations.

%%%%%%%%%%%%%%%%%%%%%%%%%%%%%%%%%%%%%%%%%%%%%%%%%%%%%%%%%%%%
\subsection{Mean Curvature Flow}
\label{subsec:mean-curvature-flow}
%%%%%%%%%%%%%%%%%%%%%%%%%%%%%%%%%%%%%%%%%%%%%%%%%%%%%%%%%%%%

Mean curvature flow evolves a hypersurface in its normal direction with
speed determined by mean curvature.

Schematically,

\[
\boxed{
\frac{\partial F}{\partial t}
=
-\vec H,
}
\]

where \(\vec H\) denotes the mean curvature vector under the chosen sign
convention.

The flow tends to reduce area.

It is therefore analogous to gradient descent for the area functional.

%%%%%%%%%%%%%%%%%%%%%%%%%%%%%%%%%%%%%%%%%%%%%%%%%%%%%%%%%%%%
\subsection{Shrinking Sphere}
\label{subsec:mean-curvature-shrinking-sphere}
%%%%%%%%%%%%%%%%%%%%%%%%%%%%%%%%%%%%%%%%%%%%%%%%%%%%%%%%%%%%

A round sphere remains round under mean curvature flow while its radius
decreases.

For an \(m\)-dimensional sphere of radius \(r(t)\) in
\(\mathbb{R}^{m+1}\), using the standard inward mean-curvature-flow
convention,

\[
\boxed{
r'(t)
=
-\frac{m}{r(t)}.
}
\]

Hence

\[
r(t)^2
=
r(0)^2-2mt.
\]

The sphere collapses to a point in finite time.

This is a basic example of singularity formation.

%%%%%%%%%%%%%%%%%%%%%%%%%%%%%%%%%%%%%%%%%%%%%%%%%%%%%%%%%%%%
\subsection{Ricci Flow}
\label{subsec:ricci-flow}
%%%%%%%%%%%%%%%%%%%%%%%%%%%%%%%%%%%%%%%%%%%%%%%%%%%%%%%%%%%%

Ricci flow evolves a Riemannian metric according to

\[
\boxed{
\frac{\partial g}{\partial t}
=
-2\operatorname{Ric}(g).
}
\]

It was introduced by Richard Hamilton.

Ricci flow can be viewed as a nonlinear diffusion equation for the
geometry itself.

Regions of curvature interact and evolve over time.

%%%%%%%%%%%%%%%%%%%%%%%%%%%%%%%%%%%%%%%%%%%%%%%%%%%%%%%%%%%%
\subsection{Scaling of Ricci Flow}
\label{subsec:ricci-flow-scaling}
%%%%%%%%%%%%%%%%%%%%%%%%%%%%%%%%%%%%%%%%%%%%%%%%%%%%%%%%%%%%

Ricci flow has a natural parabolic scaling.

If

\[
g(t)
\]

is a solution, then for \(c>0\),

\[
\boxed{
\widetilde g(t)
=
c\,g(t/c)
}
\]

is again a solution on the correspondingly rescaled time interval.

This scaling is important when studying singularities.

%%%%%%%%%%%%%%%%%%%%%%%%%%%%%%%%%%%%%%%%%%%%%%%%%%%%%%%%%%%%
\subsection{Ricci Flow and Curvature}
\label{subsec:ricci-flow-curvature}
%%%%%%%%%%%%%%%%%%%%%%%%%%%%%%%%%%%%%%%%%%%%%%%%%%%%%%%%%%%%

Under Ricci flow, curvature satisfies nonlinear parabolic evolution
equations.

Schematically,

\[
\boxed{
\frac{\partial}{\partial t}
\operatorname{Rm}
=
\Delta\operatorname{Rm}
+
\operatorname{Rm}*\operatorname{Rm},
}
\]

where

\[
\operatorname{Rm}*\operatorname{Rm}
\]

represents quadratic combinations of curvature terms.

The structure resembles

\[
\boxed{
\text{diffusion}
+
\text{nonlinear reaction}.
}
\]

%%%%%%%%%%%%%%%%%%%%%%%%%%%%%%%%%%%%%%%%%%%%%%%%%%%%%%%%%%%%
\subsection{Scalar Curvature Evolution}
\label{subsec:scalar-curvature-evolution}
%%%%%%%%%%%%%%%%%%%%%%%%%%%%%%%%%%%%%%%%%%%%%%%%%%%%%%%%%%%%

Under Ricci flow,

\[
\boxed{
\frac{\partial R}{\partial t}
=
\Delta_gR
+
2|\operatorname{Ric}|^2.
}
\]

This equation allows maximum-principle arguments to be applied to scalar
curvature.

It illustrates how geometric evolution becomes analytic evolution.

%%%%%%%%%%%%%%%%%%%%%%%%%%%%%%%%%%%%%%%%%%%%%%%%%%%%%%%%%%%%
\subsection{Singularities in Geometric Flows}
\label{subsec:geometric-flow-singularities}
%%%%%%%%%%%%%%%%%%%%%%%%%%%%%%%%%%%%%%%%%%%%%%%%%%%%%%%%%%%%

A geometric flow may develop a singularity in finite time.

Typically this means some curvature quantity becomes unbounded:

\[
\boxed{
|\operatorname{Rm}|
\to\infty.
}
\]

To study a singularity, one often rescales the geometry around points of
large curvature.

This process is called blow-up analysis.

%%%%%%%%%%%%%%%%%%%%%%%%%%%%%%%%%%%%%%%%%%%%%%%%%%%%%%%%%%%%
\subsection{Blow-Up Analysis}
\label{subsec:geometric-blow-up-analysis}
%%%%%%%%%%%%%%%%%%%%%%%%%%%%%%%%%%%%%%%%%%%%%%%%%%%%%%%%%%%%

Suppose curvature becomes large near points

\[
(x_k,t_k).
\]

Choose scaling factors

\[
\lambda_k\to\infty
\]

and magnify the geometry near these points.

One seeks a limiting object describing the local singularity.

The general philosophy is

\[
\boxed{
\text{rescale}
\longrightarrow
\text{extract a limit}
\longrightarrow
\text{classify the limit}.
}
\]

This strategy appears throughout nonlinear analysis.

%%%%%%%%%%%%%%%%%%%%%%%%%%%%%%%%%%%%%%%%%%%%%%%%%%%%%%%%%%%%
\subsection{Compactness}
\label{subsec:geometric-analysis-compactness}
%%%%%%%%%%%%%%%%%%%%%%%%%%%%%%%%%%%%%%%%%%%%%%%%%%%%%%%%%%%%

Compactness theorems allow sequences of geometric objects to have
convergent subsequences when appropriate quantities are controlled.

A typical problem considers

\[
(M_k,g_k)
\]

with uniform bounds on geometric quantities.

One then asks whether a subsequence converges to a limiting geometric
space.

The precise notion of convergence depends on the problem.

%%%%%%%%%%%%%%%%%%%%%%%%%%%%%%%%%%%%%%%%%%%%%%%%%%%%%%%%%%%%
\subsection{Concentration Phenomena}
\label{subsec:geometric-concentration}
%%%%%%%%%%%%%%%%%%%%%%%%%%%%%%%%%%%%%%%%%%%%%%%%%%%%%%%%%%%%

A bounded sequence of geometric objects may fail to converge strongly
because energy concentrates in smaller regions.

Schematically,

\[
\boxed{
\text{bounded energy}
\not\Rightarrow
\text{strong compactness}.
}
\]

Energy may concentrate near isolated points or lower-dimensional sets.

This occurs in harmonic maps, minimal surfaces, gauge theory, and
critical nonlinear PDEs.

%%%%%%%%%%%%%%%%%%%%%%%%%%%%%%%%%%%%%%%%%%%%%%%%%%%%%%%%%%%%
\subsection{Bubbling}
\label{subsec:geometric-bubbling}
%%%%%%%%%%%%%%%%%%%%%%%%%%%%%%%%%%%%%%%%%%%%%%%%%%%%%%%%%%%%

In some variational problems, concentration produces nontrivial limiting
objects after rescaling.

These are often called bubbles.

A sequence may therefore converge away from a finite set while additional
energy is carried by rescaled limiting solutions.

Bubbling is a characteristic phenomenon of critical geometric problems.

%%%%%%%%%%%%%%%%%%%%%%%%%%%%%%%%%%%%%%%%%%%%%%%%%%%%%%%%%%%%
\subsection{Regularity Theory}
\label{subsec:geometric-regularity}
%%%%%%%%%%%%%%%%%%%%%%%%%%%%%%%%%%%%%%%%%%%%%%%%%%%%%%%%%%%%

Geometric equations are often first solved in a weak sense.

One then asks whether weak solutions are actually smooth.

The typical progression is

\[
\boxed{
\text{weak solution}
\longrightarrow
\text{a priori estimates}
\longrightarrow
\text{improved regularity}
\longrightarrow
\text{smoothness}.
}
\]

Regularity theory is therefore essential to geometric analysis.

%%%%%%%%%%%%%%%%%%%%%%%%%%%%%%%%%%%%%%%%%%%%%%%%%%%%%%%%%%%%
\subsection{Weak Formulations}
\label{subsec:geometric-weak-formulations}
%%%%%%%%%%%%%%%%%%%%%%%%%%%%%%%%%%%%%%%%%%%%%%%%%%%%%%%%%%%%

For the equation

\[
-\Delta_g u=f,
\]

a weak formulation is obtained by multiplying by a test function
\(\varphi\) and integrating.

After integration by parts,

\[
\boxed{
\int_M
g(\nabla u,\nabla\varphi)
\,dV_g
=
\int_M
f\varphi
\,dV_g
}
\]

under suitable boundary assumptions.

This formulation requires fewer derivatives of \(u\) than the classical
equation.

%%%%%%%%%%%%%%%%%%%%%%%%%%%%%%%%%%%%%%%%%%%%%%%%%%%%%%%%%%%%
\subsection{Energy Estimates}
\label{subsec:geometric-energy-estimates}
%%%%%%%%%%%%%%%%%%%%%%%%%%%%%%%%%%%%%%%%%%%%%%%%%%%%%%%%%%%%

Energy estimates control solutions using integral quantities.

For example, testing

\[
-\Delta_g u=f
\]

against \(u\) gives formally

\[
\boxed{
\int_M
|\nabla u|^2
\,dV_g
=
\int_M
fu
\,dV_g.
}
\]

Inequalities such as Cauchy--Schwarz, Poincar\'e, and Sobolev inequalities
then produce bounds on \(u\).

These estimates form the analytic foundation of many existence and
compactness arguments.

%%%%%%%%%%%%%%%%%%%%%%%%%%%%%%%%%%%%%%%%%%%%%%%%%%%%%%%%%%%%
\subsection{Monotonicity Formulas}
\label{subsec:geometric-monotonicity-formulas}
%%%%%%%%%%%%%%%%%%%%%%%%%%%%%%%%%%%%%%%%%%%%%%%%%%%%%%%%%%%%

Many geometric PDEs possess quantities that are monotone with respect to
scale or time.

A monotonicity formula may have the schematic form

\[
\boxed{
\frac{d}{dr}\Phi(r)\geq0
}
\]

or

\[
\boxed{
\frac{d}{dt}\mathcal F(t)\leq0.
}
\]

Such formulas can reveal

\[
\boxed{
\text{regularity},
\quad
\text{singularity structure},
\quad
\text{rigidity},
\quad
\text{long-time behavior}.
}
\]

%%%%%%%%%%%%%%%%%%%%%%%%%%%%%%%%%%%%%%%%%%%%%%%%%%%%%%%%%%%%
\subsection{Comparison Geometry}
\label{subsec:comparison-geometry-analysis}
%%%%%%%%%%%%%%%%%%%%%%%%%%%%%%%%%%%%%%%%%%%%%%%%%%%%%%%%%%%%

Comparison geometry studies how curvature bounds control geometric
quantities.

A typical principle is

\[
\boxed{
\text{curvature bounds}
\Longrightarrow
\text{distance and volume bounds}.
}
\]

One compares a manifold with model spaces of constant curvature.

This converts local curvature assumptions into global geometric
information.

%%%%%%%%%%%%%%%%%%%%%%%%%%%%%%%%%%%%%%%%%%%%%%%%%%%%%%%%%%%%
\subsection{Ricci Curvature and Volume}
\label{subsec:ricci-volume-comparison}
%%%%%%%%%%%%%%%%%%%%%%%%%%%%%%%%%%%%%%%%%%%%%%%%%%%%%%%%%%%%

Lower bounds on Ricci curvature strongly constrain volume growth.

A central result is the Bishop--Gromov volume comparison theorem.

Roughly, if Ricci curvature is bounded below by that of an appropriate
constant-curvature model space, then normalized volume ratios satisfy a
monotonicity property.

This connects curvature directly to global volume geometry.

%%%%%%%%%%%%%%%%%%%%%%%%%%%%%%%%%%%%%%%%%%%%%%%%%%%%%%%%%%%%
\subsection{Ricci Curvature and Diameter}
\label{subsec:ricci-diameter}
%%%%%%%%%%%%%%%%%%%%%%%%%%%%%%%%%%%%%%%%%%%%%%%%%%%%%%%%%%%%

Positive Ricci curvature can force a manifold to have finite diameter.

A classical example is Myers' theorem.

Under a sufficiently positive Ricci lower bound,

\[
\boxed{
\operatorname{diam}(M)
\leq
\text{curvature-dependent constant}.
}
\]

Such results show how differential inequalities constrain global
topology and geometry.

%%%%%%%%%%%%%%%%%%%%%%%%%%%%%%%%%%%%%%%%%%%%%%%%%%%%%%%%%%%%
\subsection{Geometric Rigidity}
\label{subsec:geometric-rigidity}
%%%%%%%%%%%%%%%%%%%%%%%%%%%%%%%%%%%%%%%%%%%%%%%%%%%%%%%%%%%%

A rigidity theorem states that achieving equality in an inequality or
extremal estimate forces a geometric object to have a special form.

A typical structure is

\[
\boxed{
\text{inequality}
+
\text{equality case}
\Longrightarrow
\text{classification}.
}
\]

Rigidity is one of the recurring themes of geometric analysis.

%%%%%%%%%%%%%%%%%%%%%%%%%%%%%%%%%%%%%%%%%%%%%%%%%%%%%%%%%%%%
\subsection{Gauge Symmetry}
\label{subsec:geometric-gauge-symmetry}
%%%%%%%%%%%%%%%%%%%%%%%%%%%%%%%%%%%%%%%%%%%%%%%%%%%%%%%%%%%%

Geometric equations often possess symmetries.

For example, a geometric object may have many coordinate descriptions
representing the same underlying geometry.

This can prevent a differential equation from being strictly elliptic or
parabolic until a gauge or coordinate condition is chosen.

Thus one often uses

\[
\boxed{
\text{geometric equation}
+
\text{gauge fixing}
\longrightarrow
\text{analytic PDE}.
}
\]

This principle appears in Ricci flow, Einstein equations, and gauge
theory.

%%%%%%%%%%%%%%%%%%%%%%%%%%%%%%%%%%%%%%%%%%%%%%%%%%%%%%%%%%%%
\subsection{Yang--Mills Theory}
\label{subsec:yang-mills-geometric-analysis}
%%%%%%%%%%%%%%%%%%%%%%%%%%%%%%%%%%%%%%%%%%%%%%%%%%%%%%%%%%%%

Gauge theory provides another major source of geometric variational
problems.

For a connection \(A\) with curvature \(F_A\), the Yang--Mills energy is

\[
\boxed{
\mathcal{YM}(A)
=
\frac12
\int_M
|F_A|^2
\,dV_g.
}
\]

Critical points satisfy the Yang--Mills equation

\[
\boxed{
d_A^\ast F_A=0.
}
\]

This is a nonlinear geometric PDE with important connections to topology
and mathematical physics.

%%%%%%%%%%%%%%%%%%%%%%%%%%%%%%%%%%%%%%%%%%%%%%%%%%%%%%%%%%%%
\subsection{Geometric Analysis and Topology}
\label{subsec:geometric-analysis-topology}
%%%%%%%%%%%%%%%%%%%%%%%%%%%%%%%%%%%%%%%%%%%%%%%%%%%%%%%%%%%%

One of the deepest goals of geometric analysis is to extract topological
information from analytic and geometric structures.

The general pattern is

\[
\boxed{
\text{topology}
\longrightarrow
\text{choose geometry}
\longrightarrow
\text{solve PDE}
\longrightarrow
\text{recover topology}.
}
\]

Examples include the use of

\[
\boxed{
\text{minimal surfaces},
\quad
\text{harmonic maps},
\quad
\text{curvature equations},
\quad
\text{geometric flows}.
}
\]

%%%%%%%%%%%%%%%%%%%%%%%%%%%%%%%%%%%%%%%%%%%%%%%%%%%%%%%%%%%%
\subsection{Geometric Analysis and General Relativity}
\label{subsec:geometric-analysis-relativity}
%%%%%%%%%%%%%%%%%%%%%%%%%%%%%%%%%%%%%%%%%%%%%%%%%%%%%%%%%%%%

In general relativity, spacetime geometry is governed by Einstein's field
equations.

In geometric units, their schematic form is

\[
\boxed{
G_{\mu\nu}
+
\Lambda g_{\mu\nu}
=
8\pi T_{\mu\nu}.
}
\]

Here

\[
G_{\mu\nu}
\]

is the Einstein tensor,

\[
\Lambda
\]

is the cosmological constant, and

\[
T_{\mu\nu}
\]

is the stress-energy tensor.

The equations connect curvature with physical matter and energy.

%%%%%%%%%%%%%%%%%%%%%%%%%%%%%%%%%%%%%%%%%%%%%%%%%%%%%%%%%%%%
\subsection{Common Errors}
\label{subsec:geometric-analysis-common-errors}
%%%%%%%%%%%%%%%%%%%%%%%%%%%%%%%%%%%%%%%%%%%%%%%%%%%%%%%%%%%%

\subsubsection{Treating Geometric Operators as Euclidean Operators}

On a curved manifold,

\[
\nabla,
\qquad
\operatorname{div},
\qquad
\Delta_g
\]

depend on the metric.

Euclidean coordinate formulas cannot generally be used unchanged.

%%%%%%%%%%%%%%%%%%%%%%%%%%%%%%%%%%%%%%%%%%%%%%%%%%%%%%%%%%%%
\subsubsection{Forgetting the Volume Element}

Integration on a Riemannian manifold uses

\[
\boxed{
dV_g
=
\sqrt{\det g}\,dx^1\cdots dx^n.
}
\]

Using ordinary coordinate volume without the metric factor generally
gives the wrong integral.

%%%%%%%%%%%%%%%%%%%%%%%%%%%%%%%%%%%%%%%%%%%%%%%%%%%%%%%%%%%%
\subsubsection{Confusing Mean Curvature Conventions}

Some authors define mean curvature as the sum of principal curvatures,
while others use their average.

Signs also depend on the orientation of the normal.

Always state the convention.

%%%%%%%%%%%%%%%%%%%%%%%%%%%%%%%%%%%%%%%%%%%%%%%%%%%%%%%%%%%%
\subsubsection{Assuming Every Critical Point Is a Minimum}

The Euler--Lagrange equation identifies stationary objects.

A stationary geometric object may be

\[
\boxed{
\text{stable},
\quad
\text{unstable},
\quad
\text{a saddle point}.
}
\]

Second variation is needed to study stability.

%%%%%%%%%%%%%%%%%%%%%%%%%%%%%%%%%%%%%%%%%%%%%%%%%%%%%%%%%%%%
\subsubsection{Confusing Ricci and Scalar Curvature}

Ricci curvature is a tensor,

\[
\operatorname{Ric},
\]

while scalar curvature is its trace,

\[
R.
\]

They contain different amounts of geometric information.

%%%%%%%%%%%%%%%%%%%%%%%%%%%%%%%%%%%%%%%%%%%%%%%%%%%%%%%%%%%%
\subsubsection{Assuming Coordinates Are Geometry}

Coordinate formulas depend on the chosen chart.

The underlying geometric quantities should be coordinate invariant.

%%%%%%%%%%%%%%%%%%%%%%%%%%%%%%%%%%%%%%%%%%%%%%%%%%%%%%%%%%%%
\subsubsection{Ignoring Gauge Freedom}

A geometric PDE may appear analytically degenerate because equivalent
geometries can be represented in different coordinates or gauges.

A suitable gauge condition may be necessary.

%%%%%%%%%%%%%%%%%%%%%%%%%%%%%%%%%%%%%%%%%%%%%%%%%%%%%%%%%%%%
\subsubsection{Assuming Weak Convergence Prevents Concentration}

Weakly convergent sequences may still concentrate energy.

This is why bubbling and concentration-compactness phenomena occur.

%%%%%%%%%%%%%%%%%%%%%%%%%%%%%%%%%%%%%%%%%%%%%%%%%%%%%%%%%%%%
\subsubsection{Assuming Smooth Initial Data Prevents Singularities}

Nonlinear geometric flows can develop singularities even from smooth
initial data.

The shrinking sphere under mean curvature flow is a basic example.

%%%%%%%%%%%%%%%%%%%%%%%%%%%%%%%%%%%%%%%%%%%%%%%%%%%%%%%%%%%%
\subsubsection{Ignoring Scaling}

Geometric PDEs often possess natural scaling laws.

Scaling is essential for identifying critical quantities and analyzing
singularities.

%%%%%%%%%%%%%%%%%%%%%%%%%%%%%%%%%%%%%%%%%%%%%%%%%%%%%%%%%%%%
\subsection{Applications}
\label{subsec:geometric-analysis-applications}
%%%%%%%%%%%%%%%%%%%%%%%%%%%%%%%%%%%%%%%%%%%%%%%%%%%%%%%%%%%%

\begin{applicationbox}
\textbf{Differential Geometry.}
Analytic methods are used to construct and classify metrics, geodesics,
minimal submanifolds, and curvature structures.

\medskip

\textbf{Topology.}
Geometric PDEs can reveal global topological information and distinguish
different manifolds.

\medskip

\textbf{Partial Differential Equations.}
Geometric problems provide major classes of nonlinear elliptic and
parabolic equations.

\medskip

\textbf{Calculus of Variations.}
Minimal surfaces, harmonic maps, geodesics, and curvature equations arise
as critical points of geometric functionals.

\medskip

\textbf{General Relativity.}
Einstein's equations describe the geometry of spacetime through
curvature equations.

\medskip

\textbf{Mathematical Physics.}
Gauge theory, Yang--Mills equations, and geometric field theories combine
geometry with nonlinear analysis.

\medskip

\textbf{Image and Shape Processing.}
Curvature flows and geometric PDEs are used in smoothing, denoising,
shape evolution, and interface motion.

\medskip

\textbf{Spectral Geometry.}
Eigenvalues of geometric differential operators encode information about
shape and curvature.

\medskip

\textbf{Materials Science.}
Mean curvature and related geometric flows model surface tension and
interface evolution.

\medskip

\textbf{Optimization.}
Geometric flows can be interpreted as infinite-dimensional gradient
descent for geometric energies.
\end{applicationbox}

%%%%%%%%%%%%%%%%%%%%%%%%%%%%%%%%%%%%%%%%%%%%%%%%%%%%%%%%%%%%
\subsection{Exercises}
\label{subsec:geometric-analysis-exercises}
%%%%%%%%%%%%%%%%%%%%%%%%%%%%%%%%%%%%%%%%%%%%%%%%%%%%%%%%%%%%

\begin{exercise}[1]
Explain the basic philosophy

\[
\text{geometry}
\longleftrightarrow
\text{differential equations}.
\]
\end{exercise}

\begin{exercise}[1]
Write the local formula for the Riemannian volume element.
\end{exercise}

\begin{exercise}[1]
Define the Laplace--Beltrami operator.
\end{exercise}

\begin{exercise}[1]
Show that the Laplace--Beltrami operator reduces to the ordinary
Laplacian for the Euclidean metric.
\end{exercise}

\begin{exercise}[1]
Define a harmonic function on a Riemannian manifold.
\end{exercise}

\begin{exercise}[1]
State the Dirichlet energy of a function.
\end{exercise}

\begin{exercise}[2]
Derive the Euler--Lagrange equation for the Dirichlet energy.
\end{exercise}

\begin{exercise}[2]
Write the energy functional for a curve on a Riemannian manifold.
\end{exercise}

\begin{exercise}[2]
Derive the geodesic equation from the energy functional.
\end{exercise}

\begin{exercise}[2]
Compute the Christoffel symbols for the Euclidean metric in Cartesian
coordinates.
\end{exercise}

\begin{exercise}[2]
Verify that a plane satisfies the minimal surface equation.
\end{exercise}

\begin{exercise}[2]
Show that a constant function is harmonic.
\end{exercise}

\begin{exercise}[2]
Explain the distinction between sectional, Ricci, and scalar curvature.
\end{exercise}

\begin{exercise}[2]
Define an Einstein metric.
\end{exercise}

\begin{exercise}[2]
Explain why the Yamabe equation is nonlinear.
\end{exercise}

\begin{exercise}[2]
State the Ricci flow equation.
\end{exercise}

\begin{exercise}[2]
State the mean curvature flow equation schematically.
\end{exercise}

\begin{exercise}[3]
Compute the Laplace--Beltrami operator for a simple diagonal metric.
\end{exercise}

\begin{exercise}[3]
Derive the minimal surface equation for a graph

\[
z=u(x,y).
\]
\end{exercise}

\begin{exercise}[3]
Derive the radius equation for a sphere evolving by mean curvature flow.
\end{exercise}

\begin{exercise}[3]
Solve

\[
r'(t)=-\frac{m}{r(t)}
\]

and determine the extinction time.
\end{exercise}

\begin{exercise}[3]
Use the Bochner formula to show how nonnegative Ricci curvature
constrains harmonic functions under suitable global assumptions.
\end{exercise}

\begin{exercise}[3]
Derive the weak formulation of

\[
-\Delta_g u=f.
\]
\end{exercise}

\begin{exercise}[3]
Use the Rayleigh quotient to explain why constants correspond to the
zero eigenvalue of the Laplace--Beltrami operator on a compact connected
manifold without boundary.
\end{exercise}

\begin{exercise}[3]
Explain why concentration can prevent strong convergence even when an
energy sequence remains bounded.
\end{exercise}

\begin{exercise}[3]
Verify the parabolic scaling of Ricci flow.
\end{exercise}

\begin{exercise}[3]
Explain why curvature blow-up is a natural signal of singularity
formation.
\end{exercise}

\begin{exercise}[4]
Study the second variation of area for a minimal hypersurface and derive
the corresponding stability operator.
\end{exercise}

\begin{exercise}[4]
Derive the harmonic-map equation in local coordinates.
\end{exercise}

\begin{exercise}[4]
Study the conformal transformation law for scalar curvature and derive
the Yamabe equation.
\end{exercise}

\begin{exercise}[4]
Study the Bishop--Gromov volume comparison theorem.
\end{exercise}

\begin{exercise}[4]
Study Myers' theorem and explain how positive Ricci curvature constrains
global geometry.
\end{exercise}

\begin{exercise}[4]
Investigate a monotonicity formula for minimal surfaces or harmonic maps.
\end{exercise}

\begin{exercise}[4]
Study blow-up analysis for a simple geometric-flow singularity.
\end{exercise}

\begin{exercise}[4]
Investigate the role of gauge fixing in converting a geometric equation
into an analytically tractable PDE.
\end{exercise}

\begin{exercise}[4]
Study the Yang--Mills functional and derive its Euler--Lagrange equation.
\end{exercise}

%%%%%%%%%%%%%%%%%%%%%%%%%%%%%%%%%%%%%%%%%%%%%%%%%%%%%%%%%%%%
\subsection{Conceptual Summary}
\label{subsec:geometric-analysis-summary}
%%%%%%%%%%%%%%%%%%%%%%%%%%%%%%%%%%%%%%%%%%%%%%%%%%%%%%%%%%%%

Geometric analysis uses analytic methods to investigate geometric
structures.

Its central principle is

\[
\boxed{
\text{geometric problem}
\longrightarrow
\text{analytic equation}.
}
\]

A Riemannian metric determines geometric differential operators,
including

\[
\boxed{
\nabla,
\quad
\operatorname{div}_g,
\quad
\Delta_g.
}
\]

Harmonic functions satisfy

\[
\boxed{
\Delta_g u=0
}
\]

and arise as critical points of Dirichlet energy.

Geodesics arise as critical points of curve energy.

Minimal surfaces arise as critical points of area and satisfy

\[
\boxed{
H=0.
}
\]

Curvature appears at several levels:

\[
\boxed{
\operatorname{Rm}
\longrightarrow
\operatorname{Ric}
\longrightarrow
R.
}
\]

Curvature equations lead to nonlinear elliptic PDEs.

Geometric flows evolve geometric structures according to nonlinear
parabolic equations.

Two fundamental examples are

\[
\boxed{
\frac{\partial F}{\partial t}
=
-\vec H
}
\]

and

\[
\boxed{
\frac{\partial g}{\partial t}
=
-2\operatorname{Ric}.
}
\]

Analytic tools such as

\[
\boxed{
\text{maximum principles},
\quad
\text{Sobolev inequalities},
\quad
\text{energy estimates},
\quad
\text{compactness},
\quad
\text{blow-up analysis}
}
\]

are used to study these equations.

Geometric analysis therefore forms a major bridge connecting

\[
\boxed{
\text{analysis},
\quad
\text{geometry},
\quad
\text{topology},
\quad
\text{mathematical physics}.
}
\]

%%%%%%%%%%%%%%%%%%%%%%%%%%%%%%%%%%%%%%%%%%%%%%%%%%%%%%%%%%%%
\subsection{Section Connections}
\label{subsec:geometric-analysis-connections}
%%%%%%%%%%%%%%%%%%%%%%%%%%%%%%%%%%%%%%%%%%%%%%%%%%%%%%%%%%%%

\chapterconnection{
Geometric Analysis completes the Analysis chapter by bringing together
many of its central ideas. Real Analysis supplies convergence,
compactness, and estimates. Measure Theory and Lebesgue Integration
provide the integration framework for geometric energies. Functional
Analysis supplies weak convergence and variational methods. Operator
Theory and Spectral Theory appear through the Laplace--Beltrami operator
and its eigenvalues. Calculus of Variations produces geodesics, minimal
surfaces, harmonic maps, and curvature equations. Distribution Theory
and weak derivatives support weak formulations. Potential Theory studies
harmonic functions and Laplace equations. Approximation Theory supports
finite-dimensional methods for geometric PDEs. Asymptotic Analysis
provides scaling and blow-up techniques. Differential Geometry supplies
the underlying manifolds, metrics, connections, and curvature, while
Topology provides the global structures that geometric analysis seeks to
understand. Later chapters on Differential Equations, Optimization, and
Numerical Mathematics develop the analytic and computational methods
used to solve these geometric equations in greater detail.
}

%%%%%%%%%%%%%%%%%%%%%%%%%%%%%%%%%%%%%%%%%%%%%%%%%%%%%%%%%%%%
% SECTION SOURCE TRACKING
%%%%%%%%%%%%%%%%%%%%%%%%%%%%%%%%%%%%%%%%%%%%%%%%%%%%%%%%%%%%

%------------------------------------------------------------
% 5.18 Geometric Analysis
%------------------------------------------------------------
%
% Source basis:
%   - Standard Riemannian geometric analysis.
%   - Standard variational methods in geometry.
%   - Standard elliptic and parabolic geometric PDE theory.
%   - Standard geometric-flow theory.
%
% Used for:
%   - geometric-analysis philosophy;
%   - Riemannian metrics and volume;
%   - gradients and divergence on manifolds;
%   - Laplace--Beltrami operator;
%   - harmonic functions;
%   - Dirichlet energy;
%   - geodesics;
%   - minimal surfaces;
%   - mean curvature;
%   - first variation of area;
%   - harmonic maps;
%   - curvature tensors;
%   - sectional, Ricci, and scalar curvature;
%   - Einstein metrics;
%   - conformal geometry;
%   - Yamabe problem;
%   - elliptic geometric equations;
%   - maximum principles;
%   - Bochner formula;
%   - geometric inequalities;
%   - Sobolev and Poincare inequalities;
%   - spectral geometry;
%   - mean curvature flow;
%   - Ricci flow;
%   - singularity formation;
%   - blow-up analysis;
%   - compactness and concentration;
%   - bubbling;
%   - regularity;
%   - weak formulations;
%   - energy estimates;
%   - monotonicity formulas;
%   - comparison geometry;
%   - rigidity;
%   - gauge symmetry;
%   - Yang--Mills theory;
%   - topology and mathematical-physics connections.
%
% Notes:
%   - Differential Geometry develops manifolds, connections, curvature,
%     and Riemannian geometry from the geometric viewpoint.
%   - Calculus of Variations develops the variational principles used
%     throughout this section.
%   - Functional Analysis develops weak convergence and operator methods.
%   - Differential Equations later develops elliptic, parabolic, and
%     nonlinear PDE methods in greater generality.
%   - Topology supplies the global invariants and classification
%     questions connected to geometric PDEs.
%   - This section completes Chapter 5: Analysis.
%
%%%%%%%%%%%%%%%%%%%%%%%%%%%%%%%%%%%%%%%%%%%%%%%%%%%%%%%%%%%%